\documentclass[12pt,a4paper]{article}

\usepackage{fontspec}
\usepackage{polyglossia}

\setmainlanguage{serbian}
\setmainfont{Times New Roman}

\usepackage{graphicx}

\usepackage{amsmath}
\usepackage{amsthm}

\usepackage{xcolor}

\usepackage[a4paper,margin=2.5cm]{geometry}



\newtheorem{definition}{Дефиниција}
\newtheorem{theorem}{Теорема}
\newtheorem{lemma}{Лема}


\begin{document}



\title{\textbf{Мој лични и академски развој}}

\author{XXX}

\date{\today}

\maketitle



\tableofcontents

\newpage



\section{Увод}

Моје име је \textbf{Марија}. Студирам на факултету
\textit{Математичка}, на смеру \textit{Информатика}.

Овај кратак рад представља мој лични живот,
образовање, интересовања и планове за будућност.



\section{Основни подаци о мени}

Рођена сам у Кини.

Тренутно живим у Београду и студирам на Математички факултету.

Моје најважније интересовање је Математика.


\subsection{Образовање}

Основно образовање завршила сам у школи Владислав Рибникар.

Након тога сам похађала средњу школу Савремена гимназија.

Тренутно сам студенткиња Математички факултету.


\subsection{Интересовања}

Моја интересовања су:

\begin{itemize}
    \item Читање романе
    \item Свирање виолине
    \item Провески са подителјима
\end{itemize}

Поред тога, у слободно време волим:

\begin{enumerate}
    \item да будем сама код кући
    \item да радим задатке из математике
    \item да пешаћим са другарима
\end{enumerate}



\section{Математички пример}

Једна од познатих математичких формула је:

\[
E = mc^2
\]

Ова формула повезује енергију, масу и брзину светлости.



\section{Основни појмови}

\begin{definition}

Лични развој представља процес унапређивања
знања, способности, навика и других личних особина.

\end{definition}



\begin{theorem}

Редовно учење и вежбање доприносе развоју
знања и практичних способности.

\end{theorem}



\begin{lemma}

Добра организација времена може олакшати
извршавање студентских обавеза.

\end{lemma}



\section{Моје вештине}

\begin{table}[h]

\centering

\begin{tabular}{|c|c|c|}

\hline

\textbf{Вештина} & \textbf{Ниво} & \textbf{Напомена} \\

\hline

Рад на рачунару & Добар & Свакодневно коришћење \\

\hline

HTML и CSS & Основни & Факултетски пројекти \\

\hline

Енглески језик & Добар & Коришћење у учењу \\


\hline

\end{tabular}

\caption{Преглед мојих вештина}

\end{table}



\section{Фотографија}

На слици је приказан део мог личног интересовања.

\begin{figure}[h]

\centering

\includegraphics[width=0.45\textwidth]{slike/moja_slika.jpg}

\caption{Моја фотографија}

\end{figure}



\section{Моји циљеви}

Мој главни циљ је да успешно завршим студије.

Желим да наставим да развијам своје знање и способности.

Посебно ми је важно да стекнем практично искуство.

\textbf{Учење и напредовање су важан део мог живота.}

\textcolor{blue}{Желим да у будућности радим посао који ме занима.}



\section{Закључак}

У овом раду представио/представила сам основне информације
о себи, свом образовању, интересовањима и циљевима.

Сматрам да су образовање, труд и континуирано учење
важни за лични и професионални развој.

\end{document}
